% ============================================================================
% Fichier  : partie2/P02_C08_actifs_donnees.tex
% Titre    : Sécurité des Actifs et des Données
% Auteur   : MINKA MI NGUIDJOÏ Thierry Emmanuel
% Version  : 1.0 -- novembre 2026
% ============================================================================
% Base : brouillon « Livre III » (classification, protection des données,
% propriété). Sources : ISO/IEC 27002:2022 (mesures 5.9 à 5.14, 7.10, 8.10) ;
% NIST SP 800-88 Rev. 1 (2014). La loi n° 2024/017 est seulement mentionnée :
% ses articles sont à collationner au Journal officiel avant toute citation.
% ============================================================================

\chapter{Sécurité des Actifs et des Données}
\label{chap:actifs}
\chaptermark{Actifs et données}

\epigraph{%
    « Ce qui n'est pas inventorié n'est pas protégé. »%
}{Principe du cours}

\niveauChapitre{Fondamental}{Chapitres~\ref{chap:proprietes}
et~\ref{chap:risques}.}

\begin{boxobjectifs}
\begin{itemize}
    \item Construire un inventaire des actifs et lui associer des
          propriétaires.
    \item Définir un schéma de classification et les règles de
          manipulation qui en découlent.
    \item Protéger les données à chaque étape de leur cycle de vie.
    \item Choisir une méthode d'effacement adaptée au support.
    \item Maîtriser les supports amovibles et les fuites de données.
\end{itemize}
\end{boxobjectifs}

\bigskip

% ============================================================================
\section*{Question directrice : quelles données ont quitté LiZbiZ, et le savait-on~?}
% ============================================================================

Après l'attaque, le comité demande quelles données figuraient dans les
40~Go copiés. Personne ne sait répondre~: aucun inventaire ne dit quels
fichiers contenaient des données de clients, ni qui en était
responsable. Sans inventaire ni classification, on ne peut ni protéger,
ni mesurer une fuite, ni prévenir les personnes concernées.

\medskip
\textbf{Plan du chapitre.}
\begin{enumerate}
    \item L'inventaire et les propriétaires (section~\ref{sec:c08-inventaire}).
    \item La classification (section~\ref{sec:c08-classification}).
    \item Le cycle de vie des données (section~\ref{sec:c08-cycle}).
    \item L'effacement sûr (section~\ref{sec:c08-effacement}).
    \item Supports amovibles et fuites (section~\ref{sec:c08-fuites}).
\end{enumerate}

\separateur

% ============================================================================
\section{L'Inventaire et les Propriétaires}
\label{sec:c08-inventaire}
% ============================================================================

L'ISO/IEC~27002 fait de l'inventaire des informations et des actifs
associés la première des mesures organisationnelles sur les
actifs~\cite{iso27002_2022}. Chaque actif a un \termecle{propriétaire}~:
une personne de l'activité, et non de l'informatique, qui répond de sa
protection, fixe ses besoins et valide qui y accède.

\begin{tableaucours}{L{3.0cm}L{2.6cm}L{2.6cm}L{1.6cm}L{2.4cm}}{%
    \tableauHeader{Actif} &
    \tableauHeader{Propriétaire} &
    \tableauHeader{Emplacement} &
    \tableauHeader{Classe} &
    \tableauHeader{Supports}
}
Fichier des clients & Directeur commercial & Base de Douala & Confidentiel
& Serveur BD-01 \\
Factures & Directrice financière & Application intégrée & Interne &
Serveurs APP-01, APP-02 \\
Dossiers clients numérisés & Directeur commercial & Partage FS-02 (Garoua)
& Confidentiel & Serveur FS-02 \\
Tarifs & Directeur commercial & Messagerie, partage & Confidentiel &
Service en nuage \\
Catalogue public & Communication & Portail web & Public & Serveur WEB-01 \\
\end{tableaucours}
\vspace{-0.8em}
\begin{center}{\small\itshape Tableau~8.1. Extrait de l'inventaire des actifs
d'information de LiZbiZ (après l'attaque)}\end{center}

% ============================================================================
\section{La Classification}
\label{sec:c08-classification}
% ============================================================================

\begin{boxdef}
\termecle{Classification}~: affectation d'une information à un niveau de
sensibilité, selon l'impact de sa divulgation, de son altération ou de sa
perte. Chaque niveau entraîne des règles de \emph{marquage} et de
\emph{manipulation}.
\end{boxdef}

\begin{tableaucours}{L{2.2cm}L{3.4cm}L{6.6cm}}{%
    \tableauHeader{Niveau} &
    \tableauHeader{Impact d'une fuite} &
    \tableauHeader{Règles de manipulation}
}
Public & Aucun & Publication libre \\
Interne & Faible & Diffusion au personnel~; pas de publication \\
Confidentiel & Important & Accès sur besoin d'en connaître~; chiffrement
en transit et sur les postes nomades~; pas de stockage personnel \\
Secret & Grave & Accès nominatif validé par le propriétaire~; chiffrement
partout~; journalisation de chaque accès \\
\end{tableaucours}
\vspace{-0.8em}
\begin{center}{\small\itshape Tableau~8.2. Un schéma de classification à
quatre niveaux}\end{center}

\begin{boxalert}
\textbf{Violation typique.} Classer tout en « confidentiel ».

\smallskip
\textbf{Règle~:} une classification trop haute est ignorée, parce
qu'elle rend le travail impossible. On classe au plus juste, par
défaut au niveau « interne », et le propriétaire justifie chaque niveau
supérieur.
\end{boxalert}

\begin{boxinfo}
\textbf{Les données personnelles.} Les données qui identifient une
personne (nom, téléphone, adresse) relèvent en outre d'obligations
légales~: la loi n°~2024/017 relative à la protection des données à
caractère personnel en fixe le cadre au Cameroun. Le cours en présente
les principes au Chapitre~\ref{chap:droit}~; ses articles sont à lire
dans le texte publié au Journal officiel avant toute application.
\end{boxinfo}

% ============================================================================
\section{Le Cycle de Vie des Données}
\label{sec:c08-cycle}
% ============================================================================

\begin{tableaucours}{L{2.6cm}L{4.6cm}L{5.0cm}}{%
    \tableauHeader{Étape} &
    \tableauHeader{Risque} &
    \tableauHeader{Mesure}
}
Création & Donnée non classée & Classification à la création \\
Stockage & Accès trop larges, vol du support & Contrôle d'accès,
chiffrement \\
Usage & Copie, capture d'écran & Moindre privilège, journalisation \\
Partage & Envoi au mauvais destinataire & Canaux approuvés, chiffrement \\
Archivage & Perte de lisibilité, conservation trop longue & Durées fixées,
formats pérennes \\
Destruction & Données récupérables & Effacement sûr, certificat de
destruction \\
\end{tableaucours}
\vspace{-0.8em}
\begin{center}{\small\itshape Tableau~8.3. Les risques et les mesures du cycle
de vie}\end{center}

\begin{boxscenario}{Opération GAROUA. Des dossiers qui n'auraient pas dû être là}
\textbf{Constatation.} Le partage FS-02 contenait des dossiers de clients
numérisés depuis 2015, dont des clients partis depuis plus de cinq ans.
Aucune durée de conservation n'était fixée.

\medskip
\textbf{Analyse} (confiance élevée). Une partie des données volées
n'avait plus d'utilité pour LiZbiZ. Une règle de conservation, appliquée,
aurait réduit l'ampleur de la fuite~: la donnée la mieux protégée est
celle que l'on ne garde plus.
\end{boxscenario}

% ============================================================================
\section{L'Effacement Sûr}
\label{sec:c08-effacement}
% ============================================================================

Supprimer un fichier ne l'efface pas~: le système libère la place, mais
les données restent lisibles jusqu'à ce qu'elles soient écrasées. Le NIST
SP~800-88 distingue trois niveaux d'assainissement~\cite{nist80088}.

\begin{tableaucours}{L{2.2cm}L{5.4cm}L{4.6cm}}{%
    \tableauHeader{Niveau} &
    \tableauHeader{Principe} &
    \tableauHeader{Exemple}
}
Effacer (\textit{clear}) & Écraser par les commandes normales du
support & Réécriture complète d'un disque réutilisé en interne \\
Purger (\textit{purge}) & Rendre la récupération impraticable même en
laboratoire & Effacement cryptographique, commande d'effacement sécurisé
du disque \\
Détruire (\textit{destroy}) & Rendre le support inutilisable &
Broyage, désintégration \\
\end{tableaucours}
\vspace{-0.8em}
\begin{center}{\small\itshape Tableau~8.4. Les niveaux d'assainissement du
NIST SP~800-88}\end{center}

\begin{boxprocedure}
\textbf{Sortir un support de l'inventaire}
\begin{enumerate}
    \item Identifier le support et la classe la plus haute des données
          qu'il a contenues.
    \item Choisir le niveau d'assainissement (effacer, purger, détruire)
          selon cette classe et la destination du support.
    \item Exécuter, et vérifier par échantillonnage que les données sont
          illisibles.
    \item Établir un certificat~: support, méthode, date, opérateur,
          vérification.
    \item Mettre à jour l'inventaire.
\end{enumerate}
\end{boxprocedure}

% ============================================================================
\section{Supports Amovibles et Fuites}
\label{sec:c08-fuites}
% ============================================================================

Les supports amovibles ont permis à Stuxnet de franchir un réseau isolé
(Chapitre~\ref{chap:attaques}) et permettent à un salarié de partir avec
un fichier. Les mesures~: bloquer les ports par défaut et n'autoriser que
des supports enregistrés~; chiffrer les supports autorisés~; analyser tout
support entrant. Les outils de prévention des fuites (DLP) repèrent les
données classées qui sortent par la messagerie, le web ou un support,
à condition que la classification existe.

\begin{boxmetier}{Administrateur systèmes et réseaux}
\textbf{Tenir l'inventaire technique.} L'administrateur tient la liste des
serveurs, des postes, des logiciels et de leurs versions, rapprochée
chaque mois de ce que l'on découvre sur le réseau~: un équipement inconnu
est un signal.
\end{boxmetier}

\begin{boxmetier}{RSSI}
\textbf{Faire classer par les métiers.} Le RSSI fournit le schéma et les
règles~; les propriétaires classent. Il vérifie que les règles de
manipulation sont appliquées, par exemple par une recherche de fichiers
confidentiels hors des emplacements prévus.
\end{boxmetier}

\separateur

\begin{boxresume}
\textbf{Ce qu'il faut retenir de ce chapitre~:}
\begin{itemize}
    \item \textbf{L'inventaire} est la première mesure~; chaque actif a un
          propriétaire métier.
    \item \textbf{La classification} fixe des règles de marquage et de
          manipulation~; on classe au plus juste.
    \item \textbf{Chaque étape du cycle de vie} a ses risques~; une durée de
          conservation réduit l'ampleur des fuites.
    \item \textbf{Supprimer n'est pas effacer}~: effacer, purger ou
          détruire selon la classe et la destination.
    \item \textbf{Supports amovibles et fuites} se maîtrisent par le
          blocage, le chiffrement et la détection.
\end{itemize}
\end{boxresume}

\begin{boxinfo}
\textbf{Points de contrôle de l'élève.} Je sais~: construire une ligne
d'inventaire~; proposer un schéma à quatre niveaux et ses règles~;
associer une mesure à chaque étape du cycle de vie~; choisir entre
effacer, purger et détruire~; rédiger un certificat d'effacement.
\end{boxinfo}

\bigskip

\begin{boxexo}[title={\ding{45}\quad Exercice~8.1. Le propriétaire \niveaubase}]
Qui doit être propriétaire~: du fichier de paie~; de la base des
expéditions~; du serveur qui les héberge~? Justifiez.
\end{boxexo}

\begin{boxexo}[title={\ding{45}\quad Exercice~8.2. Classer \niveaubase}]
Classez avec le schéma du Tableau~8.2~: le catalogue des services~; la
liste des salaires~; le plan de continuité~; les mots de passe des
comptes d'administration~; une note de service sur les congés.
\end{boxexo}

\begin{boxexo}[title={\ding{45}\quad Exercice~8.3. Supprimer \niveaubase}]
Expliquez pourquoi un fichier supprimé, puis la corbeille vidée, peut
encore être récupéré.
\end{boxexo}

\begin{boxexo}[title={\ding{45}\quad Exercice~8.4. Assainir \niveaumoyen}]
Choisissez le niveau d'assainissement pour~: un portable de la
comptabilité donné à une école~; un disque de serveur de base de
données qui part en réparation chez le fabricant~; une clé USB perdue
puis retrouvée.
\end{boxexo}

\begin{boxexo}[title={\ding{45}\quad Exercice~8.5. Règles de manipulation \niveaumoyen}]
Rédigez les règles de manipulation d'un document « Confidentiel » pour
l'impression, l'envoi par courriel, le stockage sur un portable et la
destruction.
\end{boxexo}

\begin{boxexo}[title={\ding{45}\quad Exercice~8.6. Conservation \niveaumoyen}]
Proposez des durées de conservation pour~: les dossiers des clients
actifs~; ceux des clients partis~; les journaux de connexion~; les
factures. Indiquez pour chacune ce qui la fonde (besoin, obligation
légale à vérifier).
\end{boxexo}

\begin{boxexo}[title={\ding{45}\quad Exercice~8.7. Inventaire et détection \niveaumoyen}]
Un balayage du réseau de LiZbiZ découvre quatre machines absentes de
l'inventaire. Que faites-vous, dans quel ordre~?
\end{boxexo}

\begin{boxexo}[title={\ding{45}\quad Exercice~8.8. La fuite \niveaumoyen}]
Reprenez la question directrice. Quelles informations auraient permis de
dire, dès le lundi, quelles données figuraient dans les 40~Go~?
\end{boxexo}

\begin{boxexo}[title={\ding{45}\quad Exercice~8.9. Les supports \niveauexpert}]
Rédigez une politique d'usage des supports amovibles pour la BANQUE
SANAGA, et dites comment l'auditeur vérifiera son application.
\end{boxexo}

\begin{boxexo}[title={\ding{45}\quad Exercice~8.10. Un programme \niveauexpert}]
Proposez à LiZbiZ un programme de six mois pour inventorier et classer
ses informations, avec les étapes, les acteurs et les livrables.
\end{boxexo}

\bigskip

\begin{boxinfo}
\textbf{Pour aller plus loin.}
\begin{itemize}
    \item ISO/IEC~27002:2022, mesures sur les actifs~\cite{iso27002_2022}.
    \item NIST SP~800-88, assainissement des supports~\cite{nist80088}.
\end{itemize}
\end{boxinfo}

\begin{boxinfo}
\textbf{Transition.} L'attaque de Garoua a commencé par une personne
trompée. Le Chapitre~\ref{chap:humain} traite du facteur humain.
\end{boxinfo}
